\documentclass[a4paper]{article}
% Español (México) con XeLaTeX: fontspec para fuentes Unicode, polyglossia
% para la división silábica y los nombres del documento en la variante mexicana.
\usepackage{fontspec}
\usepackage{polyglossia}
\setdefaultlanguage[variant=mexican]{spanish}
% Los nombres chinos van como «pinyin (original)»: xeCJK da a los caracteres
% Han su propia fuente; sin ella XeLaTeX los omite con un aviso.
% FandolSong no cubre algunos caracteres de nombres (U+73FA); WenQuanYi Zen Hei sí.
\usepackage[AutoFallBack=true]{xeCJK}
\setCJKmainfont{FandolSong-Regular.otf}
\setCJKfallbackfamilyfont{\CJKrmdefault}{WenQuanYi Zen Hei}
% polyglossia no traduce los nombres de listings.
\AtBeginDocument{\providecommand{\lstlistingname}{}\renewcommand{\lstlistingname}{Listado}%
  \providecommand{\lstlistlistingname}{}\renewcommand{\lstlistlistingname}{Índice de listados}}
\usepackage{amsmath, amssymb}
\usepackage{graphicx}
\usepackage[margin=2.5cm]{geometry}
\usepackage[most]{tcolorbox}
\usepackage{etoolbox}
\usepackage{listings}
\usepackage{hyperref}
\usepackage{booktabs}
\usepackage{subcaption}
\usepackage{float}
\newtcolorbox{knowledgebox}[1]{enhanced, colback=blue!5!white, colframe=blue!75!black, colbacktitle=blue!75!black, coltitle=white, fonttitle=\bfseries, title={#1}, attach boxed title to top left={yshift=-2mm, xshift=2mm}, boxrule=1pt, sharp corners}
\newtcolorbox{importantbox}[1]{enhanced, colback=yellow!10!white, colframe=yellow!80!black, colbacktitle=yellow!80!black, coltitle=black, fonttitle=\bfseries, title={#1}, sharp corners}
\newtcolorbox{warningbox}[1]{enhanced, colback=red!5!white, colframe=red!75!black, colbacktitle=red!75!black, coltitle=white, fonttitle=\bfseries, title={#1}, sharp corners}
\lstset{language=Python, basicstyle=\ttfamily\small, keywordstyle=\color{blue}, stringstyle=\color{red!60!black}, commentstyle=\color{green!60!black}, breaklines=true, frame=single, numbers=left, numberstyle=\tiny\color{gray}, captionpos=b, extendedchars=true}

\newcommand{\notetitle}{Reconocimiento de patrones de Attention Head y Attention Sink}
\newcommand{\noteauthors}{Elaboradas a partir de materiales públicos del curso}
\newcommand{\notedate}{2025}
\newcommand{\videochannel}{Wudaokou Nashi (五道口纳什)}
\newcommand{\videopublishdate}{2025}
\newcommand{\videoduration}{aproximadamente 21 minutos}
\newcommand{\videourl}{}
\newcommand{\videocoverpath}{cover.jpg}
\newcommand{\repourl}{https://github.com/hqhq1025/ai-course-notes}


% Footer with repo info
\usepackage{fancyhdr}
\pagestyle{fancy}
\fancyhf{}
\fancyfoot[C]{\thepage}
\fancyfoot[R]{\footnotesize\color{gray}\href{\repourl}{AI Course Notes}}
\renewcommand{\headrulewidth}{0pt}
\renewcommand{\footrulewidth}{0pt}

\begin{document}
\begin{titlepage}
\centering
{\Large Notas del curso\par}\vspace{1.2cm}
{\huge\bfseries \notetitle\par}\vspace{0.8cm}
{\large \noteauthors\par}\vspace{0.3cm}
{\large \notedate\par}\vspace{1.2cm}
\ifdefempty{\videocoverpath}{}{\includegraphics[width=0.82\textwidth,height=0.45\textheight,keepaspectratio]{\videocoverpath}\par}
\vfill
\begin{tcolorbox}[width=0.9\textwidth, colback=black!2!white, colframe=black!60, sharp corners]
\textbf{Autor o canal del video}: \videochannel\par
\textbf{Fecha de publicación}: \videopublishdate\par
\textbf{Duración del video}: \videoduration\par
\ifdefempty{\videourl}{}{\textbf{Enlace del video}: \href{\videourl}{\nolinkurl{\videourl}}\par}\par
\textbf{Página del proyecto}: \href{\repourl}{\nolinkurl{\repourl}}\par
\end{tcolorbox}
\end{titlepage}
\tableofcontents
\newpage

\section{Introducción}

Las dos entregas anteriores profundizaron en la perspectiva geométrica de RoPE y en la derivación del decaimiento remoto. Esta entrega se centra en la \textbf{interpretabilidad de Attention}: qué extrae, procesa y elabora exactamente cada capa y cada cabeza en Multi-Layer Multi-Head Attention.

El punto de partida central es el trabajo de \textbf{Gated Attention} del equipo de Qwen, ganador del mejor artículo en NeurIPS, que intenta mitigar el fenómeno de Attention Sink dentro de Attention.

\begin{knowledgebox}{La interpretabilidad de modelos como línea de investigación}
Para quien no cuenta con recursos masivos de GPU, la interpretabilidad de modelos es una buena línea de investigación: basta con construir casos de prueba, ejecutar la inferencia hacia adelante y analizar los patrones de Attention para hacer investigación con sentido. En esencia, es hacer «Debug de un modelo de lenguaje grande».
\end{knowledgebox}

\section{Repaso básico de Attention}

\subsection{Pre-Softmax Attention Score}

Fórmula central de Attention:
$$
\text{Attention}(Q, K, V) = \text{softmax}\left(\frac{QK^T}{\sqrt{d_k}}\right) V
$$

A $QK^T / \sqrt{d_k}$, antes del Softmax, se le llama \textbf{Attention Logits} o \textbf{Attention Score}. Determina cuánta atención asigna el modelo, al procesar el token de la posición $m$, al token de la posición $n$.

\subsection{Diversidad Multi-Layer Multi-Head}

En el transformer, cada capa tiene varios Attention Head, y los distintos Head de las distintas capas aprenden patrones de atención muy diferentes entre sí. Al visualizar estos patrones se puede entender el mecanismo interno del modelo.

\subsection{Resumen de la sección}

El Attention Score es una ventana para entender el comportamiento del modelo. El mecanismo Multi-Head garantiza que los distintos Head puedan aprender distintos tipos de dependencias.

\section{Patrones típicos de los Attention Head}

\subsection{Patrones comunes de Head}

Mediante la construcción de secuencias de prueba y la visualización de las matrices de Attention de los distintos Head, se pueden identificar los siguientes patrones típicos:

\begin{importantbox}{Cinco patrones típicos de Attention}
\begin{enumerate}
  \item \textbf{Local Head}: atiende principalmente a los tokens adyacentes, con una forma diagonal
  \item \textbf{Strided Head}: atiende a los tokens con un paso fijo, capturando patrones periódicos
  \item \textbf{Global Head}: atiende de forma casi uniforme a todos los tokens
  \item \textbf{Position Head}: atiende a posiciones específicas (como el inicio de la oración o el inicio del párrafo)
  \item \textbf{Sink Head}: concentra una gran cantidad de atención en el primer token de la secuencia
\end{enumerate}
\end{importantbox}

\subsection{División funcional entre las distintas capas}

Patrón general:
\begin{itemize}
  \item \textbf{Capas superficiales} (inferiores): más Local Head, se encargan de la gramática y las combinaciones de palabras
  \item \textbf{Capas intermedias}: patrones mixtos, comienzan a establecerse asociaciones semánticas
  \item \textbf{Capas profundas} (superiores): más Global Head y Position Head, se encargan de la semántica de alto nivel
\end{itemize}

\subsection{Resumen de la sección}

Los distintos Head se diferencian de manera natural en distintas funciones. Las capas superficiales atienden a la sintaxis local y las capas profundas atienden a la semántica global, formando una estructura de abstracción jerárquica de abajo hacia arriba.

\section{El fenómeno de Attention Sink}

\subsection{¿Qué es Attention Sink?}

\begin{importantbox}{Definición de Attention Sink}
En los modelos de lenguaje autorregresivos, una gran cantidad de Attention Head asignan un peso de atención anormalmente alto al \textbf{primer token} de la secuencia (o a los primeros tokens), aunque estos tokens no sean importantes en el plano semántico. Este fenómeno se llama \textbf{Attention Sink}.
\end{importantbox}

\subsection{¿Por qué aparece Attention Sink?}

La naturaleza del Softmax exige que la suma de los pesos de atención sea 1. Cuando el modelo «no tiene mucho que atender» en ciertas posiciones, necesita «verter» la atención en algún lugar. El primer token se convierte en el «bote de basura» predeterminado:

\begin{warningbox}{Causas de Attention Sink}
\begin{itemize}
  \item La normalización forzada del Softmax $\to$ la atención debe repartirse en algún lugar
  \item El primer token es visible para todos los tokens posteriores (bajo el mask causal)
  \item Durante el entrenamiento, el modelo aprende de manera natural a verter ahí la atención «sin sentido»
  \item El primer token en sí no tiene un valor semántico especial
\end{itemize}
\end{warningbox}

\subsection{Impacto de Attention Sink}

\begin{itemize}
  \item \textbf{KV Cache}: en la inferencia en flujo continuo (sliding window), si se descarta el KV del primer token, la calidad de la salida del modelo se degrada de manera notable
  \item \textbf{Calidad del modelo}: los recursos de atención se desperdician en posiciones sin sentido
  \item \textbf{Interpretabilidad}: interfiere con el análisis de los patrones reales de atención del modelo
\end{itemize}

\subsection{Resumen de la sección}

Attention Sink es un subproducto de la normalización del Softmax: el primer token se convierte en el «receptor predeterminado» de la atención. Este fenómeno tiene un impacto importante en la gestión de KV Cache, la eficiencia del modelo y la interpretabilidad.

\section{Gated Attention: mitigación del Attention Sink}

\subsection{Idea central}

Gated Attention introduce un \textbf{mecanismo de compuerta} que permite al modelo, en ciertos Heads o posiciones, «no asignar atención», evitando así que se vea forzado a verter la atención sobre el primer token:

$$
\text{GatedAttention}(Q, K, V) = g \odot \text{softmax}\left(\frac{QK^T}{\sqrt{d_k}}\right) V
$$

donde $g$ es un vector de compuerta que puede aprenderse y que puede escalar la salida de atención de ciertas posiciones hasta acercarla a 0.

\subsection{Resumen de la sección}

Gated Attention, mediante el mecanismo de compuerta, le da al modelo «la posibilidad de no atender», y así mitiga el problema del Attention Sink que provoca la normalización forzada del Softmax.

\section{Síntesis y ampliación}

\begin{enumerate}
  \item Las distintas cabezas de Multi-Head Attention se especializan de manera natural en patrones Local, Global y Sink
  \item Attention Sink es un subproducto de la normalización de Softmax: el primer token funciona como un «basurero» de atención
  \item Gated Attention (mejor artículo de NeurIPS) mitiga este problema mediante compuertas
  \item La investigación en interpretabilidad de modelos solo requiere inferencia hacia adelante, lo cual la hace adecuada para investigadores con recursos limitados
\end{enumerate}

\subsection{Lecturas adicionales}
\begin{itemize}
  \item Qwen/mejor artículo de NeurIPS: Gated Attention
  \item «Efficient Streaming Language Models with Attention Sinks» (Xiao et al., 2023)
  \item Investigación de Anthropic sobre la clasificación de cabezas de atención
\end{itemize}

\end{document}
